\chapter{Boh}

\begin{theorem}[Diagonalization theorem for \(\mathrm{O}(3)\)]
    Let \(R \in \mathrm{O}(3) \subset \mathbb{C}(3)\). Then there is a sign \(\epsilon \in \{-1,\,+1\}\) and a phase \(q \in [0,\,\pi]\) and a basis \(Q_{\alpha}\), such that \(\alpha= -1,\,0,\,1\). Then the following holds:
    \begin{enumerate}
        \item \(v_{\alpha} \cdot v_{\beta} = \delta_{-\alpha,\beta}\),
        \item \(v_{\alpha}^* = v_{-\alpha}\),
        \item \(v_0 \times v_{\alpha} = i \alpha v_{\alpha}\),
        \item \dots
        \item \dots
    \end{enumerate}
\end{theorem}

\begin{proof}
    let \(w \in \mathbb{C}^3\) invariant under \(R\). Then we cn define its ortonormal vector \(w^{\bot} \in \mathbb{C}^3\) which is also invariant under \(R\). Then ...

    Thus \(R\) has always one real eigenvalue and two complex conjugate eigenvalues (maybe also real). The real one \(\lambda_0 = \epsilon\) while the complex ones are \(\lambda_{\pm 1} = e^{\pm i q}\). Then we can define
    \[
        \lambda_{\alpha} = \epsilon e^{-i \alpha q}, \qquad \alpha = -1,\,0,\,1.
    \]

    Thus
    \[
        R v_{\alpha} = v_{\alpha} \Lambda_{\alpha}, \quad R v_{\alpha}^* = v_{\alpha}^* \Lambda_{-\alpha}^* =e^{i \chi_\alpha} v_{-\alpha} \Lambda_{-\alpha}.
    \]
    since given (?) we have that \(v_{\alpha}^* = e^{-i \chi_\alpha} v_{-\alpha}\).
        [...]
    \[
        v_1 = \frac{1}{\sqrt{2}}(v_1 + v_{-1}), \qquad v_{2} = \frac{1}{\sqrt{2} i}(v_1 - v_{-1}), \qquad v_3 = v_0,
    \]
    and we can prove this basis to be ortonormal and to satisfy the properties of the theorem. Hence \(v_i \times v_j = \sum_{k=-1}^1 \epsilon_{ijk} v_k\), proving (3) and (4).
\end{proof}

\dots

This shows that defining these matrices in complex space give us the ability to diagonalize them, which is not possible in real space (except for special cases). This is the reason why we need to complexify the representation of \(\mathrm{O}(3)\) to get the irreducible representations of it.

\begin{proposition}[...]
    Let \(R \in \mathrm{O}(3) \subset \mathbb{C}(3)\). Then there is a sign \(\epsilon \in \{-1,\,+1\}\) and a phase \(\phi \in [0,\,\pi]\) and a unit vector \(\mathbf{n} \in \mathbb{R}^3\), \(n^2 = 1\), such that
    \[
        R = \epsilon \left[ \mathbf{n} \otimes \mathbf{n} + \cos(\phi) \left(1- \mathbf{n} \otimes \mathbf{n}\right) - \sin(\phi) \star \mathbf{n} \right].
    \]
    Every matrix \(R\) of this form is an element of \(\mathrm{O}(3)\). The opposite is also true: every element of \(\mathrm{O}(3) \subset \mathbb{C}(3)\) can be written in this form.
\end{proposition}

\begin{proof}
    \[
        R = \epsilon \left[ \mathbf{v}_0 \otimes \mathbf{v}_0 + (\cos(\phi)-i\sin(\phi)) \mathbf{v}_1 \otimes \mathbf{v}_{-1} \right],
    \]
    were now we can write
    \[
        \epsilon \left[ \mathbf{v}_0 \otimes \mathbf{v}_0 + \cos(\phi)\left( \mathbf{v}_1 \otimes \mathbf{v}_{-1} + \mathbf{v}_{-1} \otimes \mathbf{v}_1 \right) - i\sin(\phi)\left( \mathbf{v}_1 \otimes \mathbf{v}_{-1} - \mathbf{v}_{-1} \otimes \mathbf{v}_1 \right) \right],
    \]
    since ... now
    \[
        \sum_{\alpha} \mathbf{v}_{\alpha} \otimes \mathbf{v}_{-\alpha} = 1, \quad \implies \quad \mathbf{v}_1 \otimes \mathbf{v}_{-1} + \mathbf{v}_{-1} \otimes \mathbf{v}_1 = 1 - \mathbf{v}_0 \otimes \mathbf{v}_0,
    \]
    thus
    \[
        \mathbf{v}_1 \wedge \mathbf{v}_{-1} =  \mathbf{v}_1 \otimes \mathbf{v}_{-1} - \mathbf{v}_{-1} \otimes \mathbf{v}_1  = \star (\mathbf{v}_1 \times \mathbf{v}_{-1}) = \star (-i \mathbf{v}_0) = -i \star \mathbf{v}_0.
    \]
    So we get
    \[
        R = \epsilon \left[ \mathbf{v}_0 \otimes \mathbf{v}_0 + \cos(\phi) (1 - \mathbf{v}_0 \otimes \mathbf{v}_0) - i\sin(\phi) (-i \star \mathbf{v}_0) \right] = \epsilon \left[ \mathbf{v}_0 \otimes \mathbf{v}_0 + \cos(\phi) (1 - \mathbf{v}_0 \otimes \mathbf{v}_0) - \sin(\phi) \star \mathbf{v}_0\right].
    \]
    By setting \(\mathbf{n} = \mathbf{v}_0\) we get the desired result, and we can do this since \(\mathbf{v}_0 \in \mathbb{R}^3\) and \(\mathbf{v}_0^2 = 1\).
\end{proof}

We can name the previous expression as the \textit{angle-axis parametrization} of \(\mathrm{O}(3)\). The angle \(\phi\) is the angle of rotation, while the axis of rotation is given by the unit vector \(\mathbf{n}\). The sign \(\epsilon\) tells us whether we are dealing with a proper or an improper rotation. In particular, if \(\epsilon = +1\) we are dealing with a proper rotation, while if \(\epsilon = -1\) we are dealing with an improper rotation (a reflection). Thus we can write any matrices in \(\mathrm{O}(3)\) as
\[
    R_{\epsilon}(\phi,\,\mathbf{n}) = \epsilon \left[ \mathbf{n} \otimes \mathbf{n} + \cos(\phi) (1 - \mathbf{n} \otimes \mathbf{n}) - \sin(\phi) \star \mathbf{n} \right].
\]

\begin{example}[Exercise]
    Being \(A \in \mathbb{R}(3)\) and \(\mathbf{v}_i\) with \(i=1,2,3\) an ortonormal basis of \(\mathbb{R}^3\) such that
    \[
        A \mathbf{v}_i = \sum_{j=1}^3 \hat{A}_{ji} \mathbf{v}_j, \quad (\hat{A}_{ji}) \in \mathbb{R}(3).
    \]
    Thus \(\hat{A}\) is a representation matrix of \(A\) in the basis \(\mathbf{v}_i\).

    Now, being \(R \in \mathrm{O}(3)\) and \(\mathbf{v}_i\)  respecting
    \begin{enumerate}
        \item \(\mathbf{v}_3 = \mathbf{n}\),
        \item \(\mathbf{v}_i \cdot \mathbf{v}_j = \delta_{ij}\) for all \(i,j\),
        \item \(\mathbf{v}_i \times \mathbf{v}_j = \sum_{k=1}^3 \epsilon_{ijk} \mathbf{v}_k\).
    \end{enumerate}
    Then we can write
    \[
        R_{\epsilon}(\phi,\,\mathbf{n}) = \epsilon \left[ \mathbf{n} \otimes \mathbf{n} + \cos(\phi) (1 - \mathbf{n} \otimes \mathbf{n}) - \sin(\phi) \star \mathbf{n} \right].
    \]
    Then its action on the basis \(\mathbf{v}_i\) is given by
    \[
        R_{\epsilon}(\phi,\,\mathbf{n}) \mathbf{v}_i = \sum_{j=1}^3 \hat{R}_{\epsilon}(\phi,\,\mathbf{n})_{ji} \mathbf{v}_j,
    \]
    where
    \[
        \hat{R}_{\epsilon}(\phi,\,\mathbf{n}) = \epsilon \begin{pmatrix}
            \cos(\phi) & -\sin(\phi) & 0 \\
            \sin(\phi) & \cos(\phi)  & 0 \\
            0          & 0           & 1
        \end{pmatrix}.
    \]

    This can be obtained by explicitly calculating the action of \(R_{\epsilon}(\phi,\,\mathbf{n})\) on the basis \(\mathbf{v}_i\). For example, we have
    \[
        \begin{aligned}
            R_{\epsilon}(\phi,\,\mathbf{n}) \mathbf{v}_i & = \epsilon \left[ \mathbf{n} \otimes \mathbf{n} + \cos(\phi) (1 - \mathbf{n} \otimes \mathbf{n}) - \sin(\phi) \star \mathbf{n} \right] \mathbf{v}_i \\
                                                         & = \dots
        \end{aligned}
    \]

    \(\hat{A}\) depends on the choice of the basis \(\mathbf{v}_i\), however its trace and determinant do not depend on the choice of the basis. In particular, we have
    \[
        \begin{dcases}
            \Tr(\hat{A}) = \Tr(A), \\
            \det(\hat{A}) = \det(A).
        \end{dcases}
    \]
    Hence we now that \(\) and \(\).
    \[
        \begin{dcases}
            \det(R_{\epsilon}(\phi, \mathbf{n})) = \det(\hat{R}_{\epsilon}(\phi, \mathbf{n})) = \epsilon^3 = \epsilon, \\
            \Tr(R_{\epsilon}(\phi, \mathbf{n})) = \Tr(\hat{R}_{\epsilon}(\phi, \mathbf{n})) = (2 \cos \phi + 1) \epsilon.
        \end{dcases}
    \]
    This is very useful since it allows us to determine the angle \(\phi\) and the sign \(\epsilon\) (which are equal for either the matrix and its representation) from the trace and determinant of \(R\): we have all the element to write the angle-axis parametrization \(\hat{R}\) of \(R\).

    This result implies that we can write the action of any element of \(\mathrm{O}(3)\) on \(\mathbb{R}^3\) as a rotation around an axis, in the following form
    \[
        A \mathbf{x} = \sum_{i} \left(\sum_{j} \hat{A}_{ij} \mathbf{x}_j \right) \mathbf{v}_i.
    \]

    Let us write
    \[
        R_{+1}(\phi,\,\mathbf{n}) = \left[ \mathbf{v}_3 \otimes \mathbf{v}_3 + \cos(\phi) (1 - \mathbf{v}_3 \otimes \mathbf{v}_3) - \sin(\phi) \star \mathbf{v}_3 \right],
    \]
    then adopting the following cylindrical coordinates:
    \[
        \begin{dcases}
            \mathbf{x}_1 = \rho \cos \theta, \\
            \mathbf{x}_2 = \rho \sin \theta, \\
            \mathbf{x}_3 = z,
            \epsilon = +1,
        \end{dcases}
    \]
    we can write the action of the orthogonal matrix \(R_{+1}(\phi,\,\mathbf{n})\) on the vector \(\mathbf{x}\) as
    \[
        \begin{aligned}
            (R_{+1}(\phi,\,\mathbf{n}) \mathbf{x})_1 & = \rho \cos(\phi + 0), \\
            (R_{+1}(\phi,\,\mathbf{n}) \mathbf{x})_2 & = \rho \sin(\phi + 0), \\
            (R_{+1}(\phi,\,\mathbf{n}) \mathbf{x})_3 & = z,
        \end{aligned}
    \]
    showing that the action of \(R_{+1}(\phi,\,\mathbf{n})\) on any vector \(\mathbf{x} \in \mathbb{R}^3\) is a rotation of angle \(\phi\) around the axis \(\mathbf{n}\).
\end{example}

Thus we have shown that if \(R \in \mathrm{O}(3)\) then there are \(\epsilon=\pm 1\) and \(\phi \in [0,\,\pi]\) and \(\mathbf{n} \in \mathbb{R}^3\) with \(\mathbf{n}^2 = 1\) such that
\[
    R_{\epsilon}(\phi,\,\mathbf{n}) = \epsilon \left[ \mathbf{n} \otimes \mathbf{n} + \cos(\phi) (1 - \mathbf{n} \otimes \mathbf{n}) - \sin(\phi) \star \mathbf{n} \right].
\]
with \(\phi\)  and \(\mathbf{n}\) given by the trace and determinant of \(R\) through the following relations:
\[
    \begin{dcases}
        \epsilon = \det(R), \\
        \phi = \arccos\left( \frac{\Tr(R)}{2 \epsilon} - \frac{1}{2} \right).
    \end{dcases}
\]

We can also write the angle-axis parametrization of \(R \in \mathrm{SO}(3)\), the special orthogonal group, with only proper rotations, as \(\epsilon = 1\), but the relation is exactely the same:
\[
    R_{+1}(\phi,\,\mathbf{n}) = \mathbf{n} \otimes \mathbf{n} + \cos(\phi) (1 - \mathbf{n} \otimes \mathbf{n}) - \sin(\phi) \star \mathbf{n}.
\]

A special case is the space reflection with respect to the origin, which is given by \(\epsilon = -1\) and \(\phi = 0\) and in physical space is called \textit{parity transformation}; it is given by
\[
    P = R_{-1}(0, \mathbf{n}) = -1, \quad P \mathbf{x} = -\mathbf{x}.
\]
With this transformation we can relate the orthogonal group \(\mathrm{O}(3)\) to the special orthogonal group \(\mathrm{SO}(3)\) as follows:
\[
    \mathrm{O}(3) = \mathrm{SO}(3) \cup P \mathrm{SO}(3).
\]

At a certain point it becomes importabt to know if this parametrization is unique, i.e. if given \(R\) there is only one way to write it as \(R_{\epsilon}(\phi,\,\mathbf{n})\). If the parameters \(\epsilon\), \(\phi\), and \(\mathbf{n}\) are uniquely determined by \(R\), then the parametrization is unique. But since the first two parameters \(\epsilon\) and \(\phi\) are determined by the trace and determinant of \(R\), which are invariant under change of basis, we only need to check if the axis of rotation \(\mathbf{n}\) is uniquely determined by \(R\). This is not the case since:
\begin{itemize}
    \item If we have a rotation of angle \(\phi = \pi\) around an axis \(\mathbf{n}\), then we can also write it as a rotation of angle \(\pi\) around the axis \(-\mathbf{n}\): it is unique up to sign change;
          \[
              \begin{aligned}
                  R = \dots \\
                  \bs{\phi} = \pi \mathbf{n} = \pi (-\mathbf{n}).
              \end{aligned}
          \]
    \item If we consider \(\phi = 0\) then the parametrization is undetermined since any axis \(\mathbf{n}\) can be chosen, since the rotation of angle \(0\) around any axis is the identity transformation: it is totally undetermined;
          \[
              R = \dots \\
              \bs{\phi} = 0 \mathbf{n} = 0 \mathbf{n}' = \mathbf{0}.
          \]
\end{itemize}
It can be proven, however, that if \(0 < \phi < \pi\), strictly comprised, then the parametrization is unique, since in this case the axis of rotation \(\mathbf{n}\) is uniquely determined by \(R\).

Then if we have defined the phase \(\phi\) as strictly comprised between \(0\) and \(\pi\), then in the phase space, the angle vector \(\bs{\phi}\) has to have a norm equal or less than pi:
\[
    \vert \bs{\phi} \vert \leq \pi.
\]
This tells us that the topology of the phase space of \(\mathrm{SO}(3)\) is a ball of radius \(\pi\) with antipodal\footnote{Opposite with respect to the origin.} points (on the surface) identified, since the points on the surface of the ball, which have norm equal to \(\pi\), are identified with their antipodal points (the one with opposite sign), since they represent the same rotation.

Now the Lie algebra the special orthogonal group \(\mathrm{SO}(3)\) is
\[
    \sigma(3) = \mathfrak{so}(3) = \{ A \in \mathbb{R}(3) : A^T = -A \},
\]
then we can state the following proposition:

\begin{proposition}
    For every \(X \in \mathfrak{so}(3)\) there exists a unique vector \(\mathbf{x} \in \mathbb{R}^3\) such that
    \[
        X = - \star \mathbf{x}, \quad \star \mathbf{x} = \begin{pmatrix}
            0             & \mathbf{x}_3  & -\mathbf{x}_2 \\
            -\mathbf{x}_3 & 0             & \mathbf{x}_1  \\
            \mathbf{x}_2  & -\mathbf{x}_1 & 0
        \end{pmatrix}.
    \]
    Indeed we have seen that with the angle-axis parametrization we can write any element of \(\mathrm{SO}(3)\) with three parameters. We denote the map \(\mathbf{x} \mapsto -\star \mathbf{x}\) with
    \[
        X = \ell(\mathbf{x}),
    \]
    with the following properties:
    \begin{enumerate}
        \item \(\ell(\mathbf{x}) \ell(\mathbf{y}) = \ell(\mathbf{x} \times \mathbf{y})\),
        \item \([\ell(\mathbf{x}),\,\ell(\mathbf{y})] \mathbf{z} = \star \mathbf{x} \times \mathbf{y} \mathbf{z} - \star \mathbf{y} \times \mathbf{x} \mathbf{z}\dots  \)
    \end{enumerate}
\end{proposition}

Indeed we know that the space of antisymmetric matrices is three-dimensional, and the map \(\mathbf{x} \mapsto \star \mathbf{x}\) is a linear map from \(\mathbb{R}^3\) to \(\mathfrak{so}(3)\). Since both spaces have the same dimension, it is enough to show that this map is injective, i.e. that if \(\star \mathbf{x} = 0\) then \(\mathbf{x} = 0\). But this is true since the only vector that has zero cross product with all the other vectors is the zero vector itself. Hence the map is injective and thus it is an isomorphism, proving the proposition.

    [...]
then with the properties of the map \(\ell\) we can show that the Lie bracket of \(\mathfrak{so}(3)\) is given by
\[
    \left[\ell_i,\,\ell_j\right] = \sum_{k=1}^3 \epsilon_{ijk} \ell_k.
\]

Another useful property comes from the exponential map, since it is naturally connected to the angle-axis parametrization. Indeed we have \(\Phi \in \mathbb{R}\), \(\mathbf{n} \in S^2\) (the unit sphere in \(\mathbb{R}^3\)), then the exponential map is given by
\[
    \exp(\ell(\phi\mathbf{n})) = R_{+1}(\Phi, \mathbf{n}).
\]

\begin{proof}
    To show this, we can write the exponential map as a power series:
    \[
        \exp(\ell(\phi\mathbf{n})) = \sum_{k=0}^{\infty} \frac{1}{k!} \ell(\phi\mathbf{n})^k = \sum_{k=0}^{\infty} \frac{1}{k!} (-\phi \star \mathbf{n})^k,
    \]
    where now if we divide the odd and even powers of \(k\) we get
    \[
        \exp(-\phi \star \mathbf{n}) = \sum_{k=0}^{\infty} \frac{1}{(2k)!} (-\phi)^k (\star \mathbf{n})^{k} = 1 + \sum_{k=1}^{\infty} \frac{\phi^{2k}}{(2k)!} (\star \mathbf{n})^{2k} - \sum_{k=0}^{\infty} \frac{\phi^{2k+1}}{(2k+1)!} (\star \mathbf{n})^{2k+1}.
    \]
    Now let us compute \((\star \mathbf{n})^2\):
    \[
        (\star \mathbf{n})^2 \mathbf{x} = \dots = -1 + \mathbf{n} \otimes \mathbf{n},
    \]
    which is a projector on the plane orthogonal to \(\mathbf{n}\). Hence, since all the projectors are idempotent, we have
    \[
        \left[ (\star \mathbf{n})^2 \right]^2 = (\star \mathbf{n})^2,
    \]
    then, with the idea to compute the expansion of the exponential map, we can write
    \[
        (\star \mathbf{n})^{2k + 1}\mathbf{x} = \dots = (-1)^k \star \mathbf{n} \mathbf{x}, \quad (\star \mathbf{n})^{2k}\mathbf{x} = \dots = (-1)^k (1 - \mathbf{n} \otimes \mathbf{n}) \mathbf{x}.
    \]
    Hence the expansion takes the following form:
    \[
        \begin{aligned}
            \exp(-\phi \star \mathbf{n}) & = 1 + \sum_{k=1}^{\infty} \frac{\phi^{2k}}{(2k)!} (-1)^k (1 - \mathbf{n} \otimes \mathbf{n}) - \sum_{k=0}^{\infty} \frac{\phi^{2k+1}}{(2k+1)!} (-1)^k \star \mathbf{n} \\
                                         & = 1 + (\cos(\phi) - 1) (1 - \mathbf{n} \otimes \mathbf{n}) - \sin(\phi) \star \mathbf{n} = R_{+1}(\Phi,\,\mathbf{n}),
        \end{aligned}
    \]
    which is exactely the angle-axis parametrization of a proper rotation of angle \(\phi\) around the axis \(\mathbf{n}\), proving the desired result.
\end{proof}